\section{Further data and discussion}\label{sec:furtherdata}

\begin{simbox}
The re-measurement described in Section~\ref{sec:further} \TBM{has / has not yet} been carried out. The tables and figures in Sections~\ref{sec:simmethods} and \ref{sec:simextrap} marked \SIM{} were produced by passing \emph{synthetic} records through exactly the analysis pipeline that will be applied to the real videos. The synthetic records were generated by numerically integrating the full (non-linearised) equations of motion of Section~\ref{sec:simulation} with identical springs, the exact point-dipole force, the decay time, frame rate and noise level of the real recordings, and the coupling constant fitted to the $25\un{mm}$ point. Because the ``true'' frequencies are known exactly, these results show what each analysis method \emph{does} to a system that obeys the dipole law, which is precisely what is needed to interpret the real data. Every simulated number must be replaced by its measured counterpart; the text has been written so that only the numbers change. Section~\ref{sec:nonlinfit}, in contrast, uses the \emph{real} data of Table~\ref{tab:rawfreq}.
\end{simbox}

\subsection{Comparison of analysis methods}\label{sec:simmethods}

Table~\ref{tab:simmethods} gives the gradient of the ln--ln plot obtained by each analysis method, both over the original range $21$--$35\un{mm}$ and over all ten separations, when the underlying system has an exponent of exactly $-5$. Figure~\ref{fig:simlnln} shows the corresponding points, and Table~\ref{tab:simdata} in Appendix~\ref{app:rawdata} lists every frequency.

\begin{table}[htbp]
\centering\small
\caption{\SIM{} -- gradient of $\ln(f_2^2-f_1^2)$ against $\ln d$ recovered by each analysis method from synthetic records of a system with true exponent $-5$ ($240\un{fps}$, $\tau=25\un{s}$, $0.1\un{mm}$ rms tracking noise). In brackets: number of separations at which two frequencies were obtained.}
\label{tab:simmethods}
\setlength{\tabcolsep}{4pt}
\begin{tabular}{@{}p{8.0cm}p{2.6cm}p{3.6cm}@{}}
\toprule
Method & 21--35 mm & all 10 separations (21--60 mm)\\
\midrule
\tabinput{tables/tab_simmethods}
\bottomrule
\end{tabular}
\end{table}

\begin{figure}[htbp]
\centering
\includegraphics[width=0.9\textwidth]{F_sim_lnln}
\caption{\SIM{} -- ln--ln plot of the frequencies recovered by four analysis methods from records of a system that obeys the dipole law exactly (black line). Grey crosses: the original procedure ($T=30\un{s}$, $A=10\un{mm}$, FFT of $x_1$); red squares: $T=90\un{s}$, $A=3\un{mm}$, FFT of $x_1$; blue triangles: FFT of the normal coordinates; green circles: two-tone time-domain fit.}
\label{fig:simlnln}
\end{figure}

The table separates the contributions of the individual remedies:
\begin{itemize}
\item The \emph{original procedure} applied to an ideal dipole system returns $\nsimGradoldrect\pm\nsimGradErroldrect$ over $21$--$35\un{mm}$, not $-5$. About $0.5$ of the shortfall is therefore produced by the analysis alone (peak pulling with $T=30\un{s}$, partly offset by the hardening at $A=10\un{mm}$). Including the points beyond $35\un{mm}$, where the merged peaks return meaningless pairs, destroys the gradient completely ($\nsimGradAlloldrect$).
\item Lengthening the record to $90\un{s}$ and reducing the amplitude to $3\un{mm}$, still reading two peaks from $x_1$, gives $\nsimGradrect\pm\nsimGradErrrect$ with the rectangular window and $\nsimGradhann\pm\nsimGradErrhann$ with Hann over $21$--$35\un{mm}$; but beyond $40\un{mm}$ the single-coordinate spectra still fail, with either window.
\item The FFT of the \emph{normal coordinates} returns $\nsimGradnormal\pm\nsimGradErrnormal$ over the full range $21$--$60\un{mm}$: all ten points, no pulling. The spectra of $\eta$ and $\xi$ at $45\un{mm}$ are shown in Figure~\ref{fig:normalcoords}: where the spectrum of $x_1$ alone shows one broad peak, $\eta$ and $\xi$ each show one sharp peak at exactly the right place, $8\un{mHz}$ apart.
\item The \emph{two-tone fit} of $x_1(t)$ also recovers all ten points and the correct gradient over $21$--$35\un{mm}$; over the full range it gives $\nsimGradAllfit\pm\nsimGradAllErrfit$, slightly low because at $50$--$60\un{mm}$ the two sinusoids are so nearly degenerate that the fit's statistical uncertainties grow to $0.3$--$2.5\un{mHz}$, which is a sizeable fraction of a $2$--$5\un{mHz}$ splitting (Figure~\ref{fig:twotone} shows the fit at $35\un{mm}$, where it is essentially perfect).
\end{itemize}

\begin{figure}[htbp]
\centering
\includegraphics[width=\textwidth]{F_normal_coords}
\caption{\SIM{} -- spectra at $d=45\un{mm}$, $T=90\un{s}$, where $f_2-f_1=8\un{mHz}$ is less than one bin ($11\un{mHz}$). Left: the FFT of $x_1$ alone cannot separate the modes. Centre and right: the normal coordinates $\eta=x_1+x_2$ and $\xi=x_1-x_2$ each contain a single mode and give $f_1$ and $f_2$ directly (dotted lines: true values).}
\label{fig:normalcoords}
\end{figure}

\begin{figure}[htbp]
\centering
\includegraphics[width=\textwidth]{F_twotone_fit}
\caption{\SIM{} -- least-squares fit of two damped sinusoids (equation~\eqref{eq:twotone}) to a synthetic $90\un{s}$ record at $d=35\un{mm}$. The residuals (bottom) are pure tracking noise, and the fitted frequencies agree with the true values to $0.2\un{mHz}$.}
\label{fig:twotone}
\end{figure}

The practical conclusion is that with normal-coordinate tracking (or time-domain fitting) the FFT resolution ceases to limit the experiment, and the whole range $21$--$60\un{mm}$ becomes usable. In the real re-measurement the corresponding table will show, in addition, how much of the residual difference from $-5$ is \emph{physical} (finite size, detuning), because those effects are absent from the simulation.

\subsection{Amplitude extrapolation}\label{sec:simextrap}

Figure~\ref{fig:extrapolation} shows the segment analysis of item~5 of Section~\ref{sec:further} for an $A=10\un{mm}$ record at $d=21\un{mm}$. The anti-phase frequency of the first $15\un{s}$ segment is $\nextrapFirst\un{Hz}$; as the amplitude decays the frequency falls, linearly in $a^2$ as equation~\eqref{eq:LL} predicts, and the intercept, $\nextrapFtwo\un{Hz}$, agrees with the small-amplitude value $\nextrapTrue\un{Hz}$ to $0.02\,\%$. Applied to the real $A=10\un{mm}$ records, this procedure gives $f_2(a\to0)$ at every $d$ and, from the gradient of each line, a direct test of the coefficient $15g/8-125g^2/48$ of the non-linear model.

\begin{figure}[htbp]
\centering
\includegraphics[width=0.8\textwidth]{F_extrapolation}
\caption{\SIM{} -- anti-phase frequency of successive $15\un{s}$ segments of a $120\un{s}$, $A=10\un{mm}$ record at $d=21\un{mm}$, plotted against the squared amplitude of the segment. The linear extrapolation to $a^2=0$ recovers the small-amplitude frequency.}
\label{fig:extrapolation}
\end{figure}

\subsection{Model fits to the measured data}\label{sec:nonlinfit}

Finally, the five physical models of item~6 were fitted by weighted non-linear least squares to the twelve resolved frequencies actually measured (Table~\ref{tab:rawfreq}, $21$--$35\un{mm}$), each frequency weighted by $1/(0.005\un{Hz})^2$. The residual sums of squares are compared in Table~\ref{tab:modelfits}. Because the $\pm5\un{mHz}$ uncertainties are half-ranges rather than standard deviations, the absolute values of $\chi^2/\text{dof}$ should not be over-interpreted; the \emph{ratios} between models are meaningful.

\begin{table}[htbp]
\centering\small
\caption{Weighted least-squares fits of five models to the twelve measured frequencies at $21$--$35\un{mm}$. $p$ = number of fitted parameters; dof $=12-p$.}
\label{tab:modelfits}
\begin{tabular}{@{}p{6.2cm}cccp{4.6cm}@{}}
\toprule
Model & $p$ & $\chi^2$ & $\chi^2/\text{dof}$ & Fitted parameters\\
\midrule
\tabinput{tables/tab_modelfits}
\bottomrule
\end{tabular}
\end{table}

Three things stand out. (i) The identical-spring point-dipole model (a) is decisively the worst, with a $\chi^2/\text{dof}$ of $\nchireda$. (ii) Letting the exponent float (b) reduces $\chi^2$ six-fold and gives $n=\nnJoint\pm\nnJointErr$, in agreement with the graphical value; this confirms that the discrepancy is not an artefact of the logarithmic transformation. (iii) The finite-size model (d), which keeps the dipole physics but replaces the point magnets by uniformly magnetised cylinders, fits \emph{as well as} the free-exponent model with the same number of parameters, and the diameter it requires, $D=\nDfit\pm\nDfitErr\un{mm}$, \TBM{is / is not} consistent with the actual magnet (\TBM{$D=\ldots$}). Adding detuning on top of finite size (e) brings no further improvement ($\chi^2$ from $\nchid$ to $\nchie$ for one extra parameter) and drives the detuning down to $f_b-f_a\approx0.02\un{Hz}$; detuning alone (c) fits worse than finite size alone. The measured data therefore favour the explanation ``the magnets are not point dipoles at these separations'' over ``the springs are different'', although only the direct measurement of $f_a$, $f_b$ and $D$ can settle the matter.

The picture that emerges for the original result of $-3.75$ is: roughly $0.5$ of the shortfall from $-5$ is produced by the measurement itself (peak pulling at $T=30\un{s}$, net of the hardening at $A=10\un{mm}$; Table~\ref{tab:simmethods}); most of the remainder is consistent with the finite size of the magnets, which makes the true stiffness fall more slowly than $d^{-5}$ between $2D$ and $4D$; spring detuning may contribute at the large-$d$ end. The dipole law itself is not contradicted: it is the $d\gg D$ limit of the finite-size model, and the fitted finite-size force approaches $d^{-4}$ (gradient $-5$) for $d\gtrsim4D$. Whether the improved measurement returns an exponent of $-5.0\pm0.1$ over the extended range, or a residual discrepancy remains, is the decisive test, and it is one that the improved method is capable of making: the simulated normal-coordinate analysis recovers the exponent to $\pm0.01$.
